\documentclass[12pt, titlepage]{article}

\usepackage{booktabs}
\usepackage{tabularx}
\usepackage{longtable}
\usepackage{array}
\usepackage{hyperref}
\hypersetup{
    colorlinks,
    citecolor=blue,
    filecolor=black,
    linkcolor=red,
    urlcolor=blue
}
\usepackage{natbib}
\usepackage{float}

\input{../Comments}
\input{../Common}

\begin{document}

\title{System Verification and Validation Plan for \progname{}} 
\author{\authname}
\date{\today}
	
\maketitle

\pagenumbering{roman}

\section*{Revision History}

\begin{tabularx}{\textwidth}{p{3cm}p{2cm}X}
\toprule {\bf Name(s)} & {\bf Date} & {\bf Notes}\\
\midrule
Maham, Cieran and Eman & Oct 27, 2025 & Added initial first draft of the VnV Plan\\
Date 2 & 1.1 & Notes\\
\bottomrule
\end{tabularx}



\newpage

\tableofcontents

\listoftables

\newpage

\section{Symbols, Abbreviations, and Acronyms}

\renewcommand{\arraystretch}{1.2}
\begin{tabular}{l l} 
  \toprule		
  \textbf{symbol} & \textbf{description}\\
  \midrule 
  T & Test\\
  \bottomrule
\end{tabular}\\

Refer to Section 4 in the \href{https://github.com/M9Huynh/technically-functional/blob/32182cfc56c602466e270b7fcbad01d892911ea1/docs/SRS/SRS.pdf}{SRS document} for a table of many more symbols, abbreviations, acronyms, and definitions.

\newpage

\pagenumbering{arabic}

This document aims to outline the verification and validation that the Physiotherapy Companion app will go through during development. The document begins with a brief description as to what the software is, what it is meant to achieve, and some relevant information before explaining in detail the standards and metrics that will be used to evaluate the system.

\section{General Information}

\subsection{Summary}

The software focused on in this document is the Physiotherapy Companion application, built for smart devices with camera access to assess a patient undergoing physiotherapy on their form during assigned physiotherapy exercises.

The application will provide support and guidance for the exercise and make real-time suggestions to improve the user’s form using a video-feed. It will also provide overviews of previous exercises to allow users to see progress and consistency, assisting in forming healthy habits towards recovery and reinforce confidence.

The application will also support a more administrative function for physiotherapists, allowing them to view exercise summaries for patients assigned to them. Additionally, physiotherapists may leave comments on summaries to ensure consistent and timely usage of the application for healthy exercising.


\subsection{Objectives}

The objectives of the testing expanded upon within this document focus on general system usability, the security of user profiles and video-feed data, and the verification of stakeholder requirements.

Usage of the system should be easily explainable for first-time users with simple dialogue popups and limited visible guiding from one function to the next. Additionally, for the video-feed portion of the application’s usage, the user should be reasonably guided through the exercise and process to set them up for successful exercise execution and form recommendation acceptance.

User profile and video-feed data should be securely stored, whether locally or within a remote database accessed by the system. User data security has become an increasingly sensitive topic within the software industry of late and users may question how their data is being stored and utilized. Users should be aware of the usage of their data and the system should ensure unauthorized access outside of the communicated plan is never possible.

Meeting stakeholder requirements, whether they be data restrictive or process flow in nature, is what will validate the core features of the system. If users do not find value in the expected usage expected from the system, then, ultimately, the system has failed to perform its primary goal. Ensuring stakeholder requirements and feedback are incorporated into the system as it is developed, both initially and onward, will ensure the system is both wanted and usable to the intended audience.

Some out-of-scope objectives include performance and user-interface quality. In this sense, quality means the overall refinement and how polished the interface looks. The implementation should not inhibit usability, but the design need not be too meticulous and visually pleasing.

For performance, the system should respond in a timely manner that does not frustrate a user’s ability to perform the main activities of the system. Designing an optimized system which may use minimal amounts of resources is not one of the current objectives of the system. These things could both be expanded into and developed, but for now are considered out of scope.

\subsection{Extras}

The extras set out for the system consist of a Design Thinking Document and a User Instructional Video.

The Design Thinking Document aims to explain the process behind, and work done towards, the final design of the system. This document explains the overall flow of processing through the system and the interfacing elements, external libraries utilized, and more. The document flow is iterative, as is the design process, and aims to communicate the entire design process of the system.

The User Instructional Video aims to communicate the method of usage of the system to the primary stakeholders. The video will demonstrate the usage of the dashboard and analytics display, the video-feed  form recommendation function, as well as an explanation of the overall feedback provided for each completed exercise. AThe instructions will be patient-centred. For the physiotherapist view the patient list, comment features, and analytical display function will be explained.


\subsection{Relevant Documentation}
\begin{itemize}
  \item Problem Statement and Goals: \citet{PS}
  \item Development Plan: \citet{DevPlan}
  \item System Requirements Specification: \citet{SRS}
  \item Hazard Analysis: \citet{HA}
  \item User Instructional Video: \citet{UserVideo}
  \item Design Thinking Document: \citet{DesignThinking}
\end{itemize}
All of the above document are important artifacts for the implementation of the system. These documents outline how the implementation will be approached and what the implementation will  follow.

More documents to come here for external libraries such as OpenCV, pytest, Python Extension, and more when it is determined which libraries and external software will be used.

\section{Plan}

This section will focus on the methods to verify and validate artifacts of the project. The methods for each facet of the system will be listed below, and additional information including possible questionnaire questions can be found in the appendix at the end of this document. The section will start primarily with verification and move to validation at the end.

\subsection{Verification and Validation Team}
See \textit{Table 1: Verification and Validation Team}.
\begin{table}
  \centering
  \begin{tabular}{|c|c|m{6cm}|}
    \hline \textbf{Role} & \textbf{Member(s)} & \textbf{Description} \\ \hline
    Test Development & Maham, Eman & Development of tests for the system and system units based on functional and non-functional requirements listed within the SRS document. \\ \hline
    Test Implementation & Cieran & Implementation of a test suite for feasibly software testable requirements, developed based on the tests made by the test development team. Also assurance of software coverage by the suite. \\ \hline
    Test Execution & Matthew & Running and verification of the test suite. \\ \hline
    Requirements Verification & Vaisnavi & Verification and tracking of requirements to ensure each has been met and tested effectively. \\ \hline
    Stakeholder Validation & All & Validation with stakeholders and external users to ensure the system developed performs as users expect; focus on usability. \\ \hline
  \end{tabular}
  \caption{Verification and Validation Team}
\end{table}
\subsection{SRS Verification}

To verify the SRS, several approaches will be used across different sections of the document. In general, the document will receive verification through feedback from 4G06 teaching assistants, peer-review feedback from another design team, and a check in with a project advisor to ensure the document represents the correct motivation for the project. This check-in will be a self review by the advisor followed by a scheduled meeting to present any probable modifications to the requirements. Based on feedback from the advisor, issues will be created and handled.

To verify the functional requirements within the SRS, the above will be used in addition to software checks for functionality and a requirements checklist which will be updated upon checking functional requirements within the implementation.
For verification of non-functional requirements within the SRS, in addition to the general checks, there will be a questionnaire of the software’s usage and another checklist with a more flexible scaling system for the NFRs.

\subsection{Design Verification}

Verification of the design will rely on feedback from 4G06 teaching assistants, peer-review feedback provided by colleagues on another design team, as well as advisor and stakeholder inquiries. The feedback from teaching assistants and peer-review feedback will be created as issues within our repository and handled accordingly. For advisor and stakeholder inquiry, there will be software handouts along with a questionnaire which is expanded on more at the end of this section.

\subsection{Verification and Validation Plan Verification}

This document will also be reviewed by the teaching assistants and provided peer-review by another design team, creating issues within the GitHub repository to be dealt with accordingly.

In addition, the test suite derived from this document will be verified via code coverage and mutation testing on the software. For validation, the final outcome of each test will be discussed in the meeting with the project advisor to ensure the system is doing what it should in each test case.

\subsection{Implementation Verification}

For testing the implemented system, there will be a software testing suite developed based on the tests listed in sections 4 and 5 of this document. Additional testing by the team includes code walk-throughs of critical sections both during development and upon reaching milestones within the development cycle. There will also be static analysis on each pull request to GitHub, ensuring naming conventions and best practice are followed throughout the project.

\subsection{Automated Testing and Verification Tools}

Automated testing includes the testing suite previously mentioned in this section; derived from the tests found in this document which are based upon the requirements found within the \href{https://github.com/M9Huynh/technically-functional/blob/32182cfc56c602466e270b7fcbad01d892911ea1/docs/SRS/SRS.pdf}{SRS document} of this project. These tests will be verified themselves with mutation testing and code coverage assessments. A testing suite such as pytest or Python Extension will be utilized to implement the tests mentioned in further sections of this document. These suits contain their own internal tools for checking code coverage, and the GitHub will be using continuous integration to ensure base cases continue to pass through all pull requests as well as general compilation checks.

During development, pylint will ensure the coding standard is adhered to and continuous development within GitHub will ensure no slip-ups on merging.

\subsection{Software Validation}

For software validation, to ensure the correct system was designed to solve the problems outlined within the problem statement document, there will be system walkthroughs and interviews with stakeholders and the project advisor. These interviews will assess time to use the systems functions, understanding of the system and how it processes data, as well as a general acceptance of the system. The interview will end with questions outlined in the appendix of this document. (Section 6.2)

\section{System testing}
The tests outlined below have been divided into 3 key areas: \\
\begin{enumerate}
\item \textbf{User registration suite:} the tests listed here have been designed for the purpose of testing the system's handling of the registration process. 
\item \textbf{PT specific testing:} These tests have been tailored to test the functionalities available to the physiotherapists. 
\item \textbf{Patient focused testing:} The results from this will give insight into how well the system accomplishes its intended purpose. 
\end{enumerate}
\textbf{Note: Automatic tests will be run using a script. See Appendix for an example template of script.
The use of ‘display’ indicates both audio and visual formats are provided.}

\subsection{User registration suite} 
This section contains: user\_acc\_create, invalid\_format\_create, cannot\_verify\_creds, acc\_exists.

\begin{longtable}{|>{\raggedright\arraybackslash}p{3cm}|>{\raggedright\arraybackslash}p{12cm}|}
\caption{Software Test Cases for User Account Creation} \\
\hline
\textbf{Test Case ID} & \textbf{user\_acc\_create} \\
\hline
\endfirsthead

\hline
Control & Automatic \\
\hline
Rationale & The system must be able to create and store a new account. \\
\hline
State & User does not have an account on the application. \\
\hline
Input & Govt id (name, date of birth), license number (PT) OR invite code (patient), email, password (will be set at the end of the registration process) \\
\hline
Output & "Account created successfully" displayed on page and user is successfully added to the user database \\
\hline
How the test will be performed & 
The test will be conducted with the use of a script (see template in Appendix)  \\
\hline

\textbf{Test Case ID} & \textbf{invalid\_format\_create} \\
\hline
Control & Automatic \\
\hline
Rationale & The data input by the user has to be in the correct format for our system to accept it as valid. \\
\hline
State & User does not have an account on the application. \\
\hline
Input & Incorrect format of license number/invite code OR incorrect format of date of birth, email,  password (will be set at the end of the registration process) \\
\hline
Output & System displays 'incorrect format of one or more inputs' \\
\hline
How the test will be performed & 
The test will be conducted with the use of a script (see template in Appendix)  \\
\hline

\textbf{Test Case ID} & \textbf{cannot\_verify\_creds} \\
\hline
Control & Automatic \\
\hline
Rationale & Successful verification for both types of users is needed to use the application \\
\hline
State & User does not have an account on the application. \\
\hline
Input & Incorrect license number/mismatch of license number and govt id OR incorrect format of invite code, email, password (will be set at the end of the registration process)  \\
\hline
Output & System displays 'License number cannot verified' OR 'invalid invite code' \\
\hline
How the test will be performed & 
The test will be conducted with the use of a script (see template in Appendix)  \\
\hline

\textbf{Test Case ID} & \textbf{acc\_exists} \\
\hline
Control & Automatic \\
\hline
Rationale & New users should not already have an account on the system \\
\hline
State & Account already exists on system \\
\hline
Input & Govt id (name, date of birth), license number/invite code, email, password \\
\hline
Output & 'Account already exists. Wrong password. Try again.' \\
\hline
How the test will be performed & 
The test will be conducted with the use of a script (see template in Appendix)  \\
\hline

\end{longtable}

\subsection{PT focused testing}
This section contains: PT\_sends\_invite, PT\_adjusts\_exercise, PT\_deletes\_patient\_acc.

\begin{longtable}{|>{\raggedright\arraybackslash}p{3cm}|>{\raggedright\arraybackslash}p{12cm}|}
\caption{Software Test Cases for Physical Therapist (PT) Functions} \\
\hline

\textbf{Test Case ID} & \textbf{PT\_sends\_invite} \\
\hline
Control & Automatic \\
\hline
Rationale & A patient can only sign up if they have an invite code from the PT. \\
\hline
State & Patient does not have an account on the application while PT does. \\
\hline
Input & PT adds patient's information and contact method (where the code will be sent). \\
\hline
Output & An 'Invite sent' message is displayed on the screen \\
\hline
How the test will be performed & 
The test will be conducted with the use of a script (see template in Appendix) \\
\hline

\textbf{Test Case ID} & \textbf{PT\_adjusts\_exercise} \\
\hline
Control & Automatic \\
\hline
Rationale & The system allows the PT to obtain an overview of patient performance and adjust different parameters accordingly (eg. number of repetitions). \\
\hline
State & PT and patient have an account on application, PT has patient on their patient list \\
\hline
Input & PT accesses the patient's profile page and adds/modifies exercise \\
\hline
Output & Patient exercise criteria is modified. \\
\hline
How the test will be performed & 
The test will be conducted with the use of a script (see template in Appendix)  \\
\hline


\textbf{Test Case ID} & \textbf{PT\_deletes\_patient\_acc} \\

\hline
Control & Automatic \\
\hline
Rationale & This ensures PTs control over patients' accounts. PT can remove patient(s) profiles from the application. \\
\hline
State & PT and patient have account on application, PT has the patient on their patient list \\
\hline
Input & PT requests a specific patient's account deletion \\
\hline
Output & Account deleted from user database and PT patient list. Patient is unable to log in ('Account not found'). \\
\hline
How the test will be performed & 
The test will be conducted with the use of a script (see template in Appendix)  \\
\hline


\end{longtable}

\subsection{Patient focused testing}
This section contains: system\_exercise\_overview, system\_record\_exercise, system\_performs\_correctly, video\_store, video\_retrieval, post\_ex\_patient\_feedback,
incorrect\_form\_adjustment \\

\begin{longtable}{|>{\raggedright\arraybackslash}p{3cm}|>{\raggedright\arraybackslash}p{12cm}|}
\caption{Software Test Cases for Exercise System and Additional Functions} \\
\hline

\textbf{Test Case ID} & \textbf{system\_exercise\_overview} \\
\hline
Control & Automatic \\
\hline
Rationale & In order to record the patient performing the exercise, the system has to be able to provide instructions to the patient on how to perform it. \\
\hline
State & Patient is logged into the application and has it open on their device. Patient has selected the exercise they want to perform on their exercise list. \\
\hline
Input & Patient clicks on exercise from exercise list. \\
\hline
Output & System displays an overview of the steps required for the exercise.\\
\hline
How the test will be performed & The test will be conducted with the use of a script (see template in Appendix)  \\

\hline
\textbf{Test Case ID} & \textbf{system\_record\_exercise} \\
\hline
Control & Manual \\
\hline
Rationale & Patient performance of the exercise has to be successfully recorded to aid system analysis of its correctness. \\
\hline
State & Patient is currently logged in and has the application open on their device. Patient has selected the exercise they would like to perform from the exercise list. \\
\hline
Input & Patient clicks on the 'record' button \\
\hline
Output & Exercise successfully recorded, saved and ready for viewing \\
\hline
How the test will be performed & 
\begin{itemize}
    \setlength\itemsep{0em}
    \item Patient clicks on the 'exercise list' and selects the exercise they would like to perform
    \item System directs patient to a screen with a live video stream of themselves
    \item Patient clicks on the 'record' button
    \item System records exercise
\end{itemize} \\
\hline

\textbf{Test Case ID} & \textbf{system\_performs\_correctly} \\
\hline
Control & Manual \\
\hline
Rationale & System has to effectively perform analysis on patient exercise form to provide feedback. \\
\hline
State & Patient is recording exercise \\
\hline
Input & System receives visual of patient performing exercise \\
\hline
Output & System provides feedback to patient (e.g. instructions for the next steps in the exercise) \\
\hline
How the test will be performed & 
\begin{itemize}
    \setlength\itemsep{0em}
    \item Patient performs exercise correctly
    \item System monitors and analyzes exercise performance in real-time
    \item System provides correct feedback
\end{itemize} \\
\hline

\textbf{Test Case ID} & \textbf{video\_store} \\
\hline
Control & Automatic \\
\hline
Rationale & The system has to be able to store a recorded video of the patient performing the exercises. \\
\hline
State & Patient has performed and recorded exercise being performed. \\
\hline
Input & Patient clicks on 'Save exercise' button. \emph{Note: The team has not decided whether to store the video locally or on the user database. }\\
\hline
Output & Video has been successfully saved and is ready for viewing. \\
\hline
How the test will be performed & The test will be conducted with the use of a script (see template in Appendix)  \\
\hline

\textbf{Test Case ID} & \textbf{video\_retrieval} \\
\hline
Control & Automatic \\
\hline
Rationale & The system has to be able to retrieve the most recent recorded video in storage (either locally or database). \\
\hline
State & Patient has performed and recorded exercise being performed \\
\hline
Input & Patient OR PT clicks on video stored (local or database) for viewing. \emph{Note: The team has not decided whether to store the video locally or on the user database. }\\
\hline
Output & Video has been successfully retrieved and is ready for viewing. \\
\hline
How the test will be performed & The test will be conducted with the use of a script (see template in Appendix) \\
\hline


\textbf{Test Case ID} & \textbf{post\_ex\_patient\_feedback} \\
\hline
Control & Automatic \\
\hline
Rationale & The system has to be able to store patient feedback of the performed exercise (e.g. the level of (perceived) difficulty of the exercise) \\
\hline
State & Patient has performed and recorded exercise being performed and system has successfully performed analysis of the performance \\
\hline
Input & Patient input into feedback form (See Appendix \emph{Patient post-exercise feedback form}). \\
\hline
Output & Answers are stored successfully \\
\hline
How the test will be performed & The test will be conducted with the use of a script (see template in Appendix) \\
\hline

\textbf{Test Case ID} & \textbf{incorrect\_form\_adjustment} \\
\hline
Control & Manual \\
\hline
Rationale & The system has to be able to differentiate between correct and incorrect forms to provide accurate feedback. \\
\hline
State & Patient is performing the selected exercise \\
\hline
Input & Incorrect form during exercise performance \\
\hline
Output & Adjustment directives displayed to patient \\
\hline
How the test will be performed & 
\begin{itemize}
    \setlength\itemsep{0em}
    \item Patient performs exercise with incorrect form
    \item System detects the incorrect posture/movement
    \item System provides real-time adjustment instructions
\end{itemize} \\
\hline

\end{longtable}

\newpage

\section{Nonfunctional Requirements Evaluation}
\begin{longtable}{|>{\raggedright\arraybackslash}p{3cm}|>{\raggedright\arraybackslash}p{12cm}|}
\caption{Non-Functional Requirements Testing} \\
\hline
\textbf{Test Case ID} & \textbf{ui\_testing} \\


\hline
Type & Static \\
\hline
Rationale & Ensure user interface meets usability standards and stakeholder expectations \\
\hline
State & N/A \\
\hline
Input & User interface elements and survey responses \\
\hline
Output & Modified user interface based on feedback \\
\hline
How the test will be performed & 
\begin{itemize}
    \setlength\itemsep{0em}
    \item A survey will be given to users (see 'Stakeholders' in SRS for the type of users) to test different interface elements
    \item See Appendix for survey questions
    \item Based on the feedback received, the user interface will be modified in order to satisfy additional requirements
\end{itemize} \\
\hline

\textbf{Test Case ID} & \textbf{ex\_int\_speed} \\
\hline
Type & Automatic \\
\hline
Rationale & Ensure exercise interface loads within acceptable time limits for good user experience \\
\hline
State & User is logged in and has application open on their device \\
\hline
Input & User clicks on exercise interface \\
\hline
Output & Interface loads within 2 seconds \\
\hline
How the test will be performed & The test will use external tools (Postman) to measure API latency. In order to replicate real world conditions, 
an increasing amount of simulated user loads can be used. Statistical analysis will be performed on the collected results to provide a more accurate value. 
\\
\hline

\textbf{Test Case ID} & \textbf{fail\_recovery} \\
\hline
Type & Dynamic \\
\hline
Rationale & The system must be able to restore its pre-failure state 95\% of the time within 5 minutes \\
\hline
State & User is currently using the application on their device \\
\hline
Input & System failure (simulated) \\
\hline
Output & System restarts and recovers within 5 minutes \\
\hline
How the test will be performed & 
\begin{itemize}
    \setlength\itemsep{0em}
    \item A failure scenario will be simulated
    \item Recovery time and success rate will be measured
    \item System's state will be compared to pre-failure state to confirm if properly restored.
\end{itemize} \\
\hline

\textbf{Test Case ID} & \textbf{Interrupted\_vid} \\
\hline
Type & Dynamic \\
\hline
Rationale & The system must be able to attempt resubmission of video after restoring internet connection \\
\hline
State & A video of the user performing the exercise is being uploaded \\
\hline
Input & Internet connection interruption during upload \\
\hline
Output & Video is successfully re-queued for upload to the system. Re-upload is resumed at the point of failure \\
\hline
How the test will be performed & 
\begin{itemize}
    \setlength\itemsep{0em}
    \item Network is disconnected during video upload
    \item Connection is restored after interruption
    \item The system detects interruption and queues video for re-upload
    \item Confirm upload resumes from point of failure (not restarting)
    \item Video integrity is confirmed after successful upload
\end{itemize} \\
\hline

 & \textbf{PR-RFT 3, PR-RFT 4, PR-CR1, PR-SER 1} \\

Rationale & These requirements address advanced system capabilities beyond current testing scope \\
\hline
How the test will be performed & 
\begin{itemize}
    \setlength\itemsep{0em}
    \item These requirements are unable to be tested in a timely manner due to project constraints (time, money)
    \item Testing deferred to future development phases
    \item Requirements marked for re-evaluation in next project iteration
\end{itemize} \\
\hline

\end{longtable}
	

\subsection{Traceability Between Test Cases and Requirements}

Please refer to the following Microsoft Excel file: \href{https://mcmasteru365-my.sharepoint.com/:x:/g/personal/dieboltc_mcmaster_ca/EYBm9tg_K8JDkVQ9Z3X4dZIBOX3TODsKXdwZCJ0CxEYkqQ?e=aGQzxo}{Traceability Matrix}

\section{Unit Test Description}

\subsection{Unit Tests}

This section outlines the unit testing approach, scope, and representative test cases for each module. Detailed input/output data and executable tests are maintained in the project’s testing repository.

\subsection{Unit Testing Scope}

All major modules developed for the system are within the scope of unit testing, including authentication, organization linking, exercise management, video capture, computer vision, reporting, data storage, and user interface components.

Third-party libraries and SDKs (e.g., video encoding, encryption frameworks) are out of scope for direct unit testing. These are verified indirectly through integration and acceptance testing. Modules responsible for UI styling and documentation display are of lower priority since their risk level is minimal; they will be verified primarily through automated UI snapshot tests.

Testing focuses on functional correctness, data integrity, and compliance with fitness criteria. Each module includes both \textbf{black box} (external behavior) and \textbf{white box} (internal logic) tests.

\subsection{Tests for Functional Requirements}

Most verification for the functional requirements will occur through automated unit tests. Each module below outlines representative tests that validate the implementation of key requirements.

\subsubsection{Module 1: Authentication and Account Management}

This module verifies that only authorized users can access protected features and that account creation, login, and security rules comply with FR 1--4 and SR AR 1--4.

\textbf{Test ID: AUTH-UT-001} \\
Type: Functional, Automatic \\
Initial State: System with no active user session. \\
Input: Attempt to access the dashboard without logging in. \\
Output: Access is denied with an authentication error message (HTTP 401 or equivalent). \\
Test Case Derivation: From FR 1 --- users must log in before accessing exercises or dashboards. \\
How Test Will Be Performed: Run automated request simulation; verify the controller response and log entry. \\[6pt]

\textbf{Test ID: AUTH-UT-002} \\
Type: Functional, Automatic \\
Initial State: Empty user database. \\
Input: Attempt to register with missing required fields or duplicate credentials. \\
Output: Registration is rejected; descriptive validation errors are returned. \\
Test Case Derivation: From FR 1.1 --- all required fields must be completed, and credentials must be unique. \\
How Test Will Be Performed: Automated validation tests using mocked data submissions. \\[6pt]

\textbf{Test ID: AUTH-UT-003} \\
Type: Functional, Automatic \\
Initial State: Active logged-in session. \\
Input: 60 minutes of inactivity simulated with timers. \\
Output: Session is automatically terminated; subsequent requests are rejected. \\
Test Case Derivation: From SR-AR 3 --- system must log out inactive users. \\
How Test Will Be Performed: Use time mocking; verify session invalidation behavior. \\[10pt]

\subsubsection{Module 2: Organization Linking and Access Codes}

This module ensures physiotherapists and patients connect securely under FR 1.2--1.3 and FR 2--3.

\textbf{Test ID: ORG-UT-001} \\
Type: Functional, Automatic \\
Initial State: Unverified physiotherapist profile. \\
Input: Submit incomplete license or clinic credentials. \\
Output: Verification denied; PT access blocked until all credentials provided. \\
Test Case Derivation: From FR 1.2 --- PT access requires all four credential fields. \\
How Test Will Be Performed: Automated validation of registration form with missing inputs. \\[6pt]

\textbf{Test ID: ORG-UT-002} \\
Type: Functional, Automatic \\
Initial State: Patient account without an organization link. \\
Input: Attempt to join with invalid and valid access codes. \\
Output: Invalid codes rejected; valid code links patient to exactly one PT and clinic. \\
Test Case Derivation: From FR 1.3 and FR 3 --- patients can only join with a valid PT-issued code. \\
How Test Will Be Performed: Automated join requests tested against mock organization data. \\[10pt]

\subsubsection{Module 3: Exercise Program Management}

Covers FR 5--6, FR 10--11, FR 22, and FR 27. Ensures exercises can be created, customized, and tracked properly.

\textbf{Test ID: PROG-UT-001} \\
Type: Functional, Automatic \\
Initial State: Exercise library initialized. \\
Input: Request for exercise list by an authorized patient. \\
Output: At least 10 knee exercises with instructions and media are returned. \\
Test Case Derivation: From FR 5 --- the system must include a complete knee-exercise library. \\
How Test Will Be Performed: Automated content verification test querying the exercise repository. \\[6pt]

\textbf{Test ID: PROG-UT-002} \\
Type: Functional, Automatic \\
Initial State: Existing exercise plan. \\
Input: PT updates sets, reps, and rest time. \\
Output: Plan updated immediately; change logged with timestamp. \\
Test Case Derivation: From FR 6 --- PTs can customize patient exercises. \\
How Test Will Be Performed: Automated API test comparing new and old plan versions. \\[6pt]

\textbf{Test ID: PROG-UT-003} \\
Type: Functional, Automatic \\
Initial State: PT sets weekly limits on exercise metrics. \\
Input: Patient attempts to exceed limit. \\
Output: System prevents action and displays a limit-exceeded message. \\
Test Case Derivation: From FR 27 --- PT-defined daily/weekly limitations must be enforced. \\
How Test Will Be Performed: Simulate plan execution and validate limit enforcement. \\[10pt]

\subsubsection{Module 4: Video Capture and Review}

Tests FR 7, FR 15, FR 26, and user-environment checks.

\textbf{Test ID: CAP-UT-001} \\
Type: Functional, Automatic \\
Initial State: Device camera available. \\
Input: Begin exercise recording. \\
Output: Video captured at $\geq$720p @ 30 fps; stored successfully for PT review. \\
Test Case Derivation: From FR 7 --- video must meet minimum quality standards. \\
How Test Will Be Performed: Automated hardware simulation verifying metadata fields. \\[6pt]

\textbf{Test ID: CAP-UT-002} \\
Type: Functional, Automatic \\
Initial State: Patient ready to record in a low-light environment. \\
Input: System performs pre-capture check. \\
Output: Warning prompts user to improve lighting before recording. \\
Test Case Derivation: From UHR-EOU 3 --- app must alert user to poor conditions. \\
How Test Will Be Performed: Simulated light-level input producing prompt verification. \\[6pt]

\textbf{Test ID: CAP-UT-003} \\
Type: Functional, Automatic \\
Initial State: Captured but incomplete exercise data. \\
Input: Submit out-of-frame video. \\
Output: System rejects submission, provides retry prompt with reason. \\
Test Case Derivation: From FR 26 --- retries allowed for insufficient data. \\
How Test Will Be Performed: Mock analysis engine returning ``insufficient data'' flag. \\[10pt]

\subsubsection{Module 5: Computer Vision and Form Analysis}

Validates FR 8, FR 12--14, FR 16, and related accuracy requirements.

\textbf{Test ID: CV-UT-001} \\
Type: Functional, Automatic \\
Initial State: Analysis model loaded. \\
Input: Test video with known joint angles. \\
Output: Angles computed within $\pm$5°; range of motion within $\pm$10\%. \\
Test Case Derivation: From FR 12 and PR-PAR 1 --- measurement accuracy thresholds. \\
How Test Will Be Performed: Automated comparison with expert-measured reference data. \\[6pt]

\textbf{Test ID: CV-UT-002} \\
Type: Functional, Automatic \\
Initial State: Exercise attempt with poor form. \\
Input: Analyze deviation from model. \\
Output: ``Incorrect form'' message specifying deviation (e.g., insufficient bend) within 1 second. \\
Test Case Derivation: From FR 13--14 --- system must identify and describe form issues. \\
How Test Will Be Performed: Simulated exercise data fed to analysis engine; timing measured. \\[6pt]

\textbf{Test ID: CV-UT-003} \\
Type: Functional, Automatic \\
Initial State: Detected dangerous deviation. \\
Input: Critical form violation ($>$30° misalignment). \\
Output: Exercise halted; warning displayed. \\
Test Case Derivation: From PR-SCR 1 --- prevent injury by halting unsafe movement. \\
How Test Will Be Performed: Inject synthetic ``critical'' deviation data; verify halt condition. \\[10pt]

\subsubsection{Module 6: Reporting and Dashboards}

Confirms FR 17--20, FR 25, and FR 18 (PT visibility).

\textbf{Test ID: RPT-UT-001} \\
Type: Functional, Automatic \\
Initial State: Completed exercise session. \\
Input: Generate session report. \\
Output: Report includes completion rate, form score, pain level, and feedback. \\
Test Case Derivation: From FR 17 --- PTs receive session reports automatically. \\
How Test Will Be Performed: Automated generation and field-validation of report content. \\[6pt]

\textbf{Test ID: RPT-UT-002} \\
Type: Functional, Automatic \\
Initial State: Patient repeatedly fails exercise. \\
Input: Record 3 failed attempts. \\
Output: Alert appears in PT dashboard. \\
Test Case Derivation: From FR 20 --- alerts for repeated incorrect attempts. \\
How Test Will Be Performed: Simulate consecutive failed submissions and verify dashboard flag. \\[10pt]

\subsubsection{Module 7: Data Security and Privacy}

Ensures data protection as required by FR 21--23 and SR IR 1--SR PR 2.

\textbf{Test ID: DATA-UT-001} \\
Type: Functional, Automatic \\
Initial State: Database ready for new patient data. \\
Input: Save patient profile with personal and medical details. \\
Output: Data stored in encrypted form; readable only via authorized session. \\
Test Case Derivation: From FR 21 and SR-IR 1 --- encryption at rest and in transit. \\
How Test Will Be Performed: Automated write-and-read test confirming ciphertext storage. \\[6pt]

\textbf{Test ID: DATA-UT-002} \\
Type: Functional, Automatic \\
Initial State: Two patient accounts exist. \\
Input: Patient A attempts to access Patient B’s records. \\
Output: Access denied; system logs unauthorized attempt. \\
Test Case Derivation: From FR 23 and SR-PR 1 --- patients may access only their own data. \\
How Test Will Be Performed: Automated access attempt verifying access-control policies. \\[10pt]

\subsection{Tests for Nonfunctional Requirements}

These tests focus on performance, robustness, scalability, and maintainability characteristics that affect the overall reliability of the system.

\subsubsection{Module 8: Performance and Latency}

\textbf{Test ID: PERF-UT-001} \\
Type: Dynamic, Automatic \\
Initial State: Application idle. \\
Input/Condition: Load exercise interface. \\
Output/Result: Interface loads within 2 seconds. \\
How Test Will Be Performed: Automated timing test using performance hooks. \\[6pt]

\textbf{Test ID: PERF-UT-002} \\
Type: Dynamic, Automatic \\
Initial State: Running video analysis. \\
Input/Condition: Introduce form deviation. \\
Output/Result: Corrective feedback appears within 1 second. \\
How Test Will Be Performed: Mock time-stamped analysis results; measure delay to user prompt. \\[10pt]

\subsubsection{Module 9: Reliability and Offline Operation}

\textbf{Test ID: REL-UT-001} \\
Type: Functional, Automatic \\
Initial State: Network unavailable. \\
Input/Condition: Attempt to upload exercise video. \\
Output/Result: Upload queued locally; auto-retries every 30 seconds until success. \\
How Test Will Be Performed: Network simulation tests queue persistence and retry logic. \\[6pt]

\textbf{Test ID: REL-UT-002} \\
Type: Functional, Automatic \\
Initial State: Device offline. \\
Input/Condition: User opens exercise list. \\
Output/Result: Cached exercises and instructions load successfully; sync resumes when online. \\
How Test Will Be Performed: Offline simulation verifying cache and synchronization behavior. \\[10pt]

\subsubsection{Module 10: Security and Backup}

\textbf{Test ID: SEC-UT-001} \\
Type: Functional, Automatic \\
Initial State: System with active data records. \\
Input/Condition: Trigger daily backup job. \\
Output/Result: Backup created and timestamped successfully. \\
How Test Will Be Performed: Automated scheduler simulation confirming backup creation. \\[6pt]

\textbf{Test ID: SEC-UT-002} \\
Type: Functional, Automatic \\
Initial State: Idle session. \\
Input/Condition: User inactive for 60 minutes. \\
Output/Result: Session automatically logs out and prompts for re-authentication. \\
How Test Will Be Performed: Time-simulation test confirming session timeout. \\[10pt]

\subsubsection{Module 11: Scalability and Extensibility}

\textbf{Test ID: EXT-UT-001} \\
Type: Functional, Automatic \\
Initial State: Base application loaded with knee module. \\
Input/Condition: Add new ``shoulder'' exercise definitions. \\
Output/Result: System loads new joint without modifying core logic. \\
How Test Will Be Performed: Add-on simulation validating modular design extensibility. \\[10pt]

\subsection{Traceability Between Test Cases and Modules}

\begin{table}[H]
	\centering
	\begin{tabular}{p{4cm} p{7cm} p{4cm}}
		\toprule
		\textbf{Module} & \textbf{Primary Requirements Covered} & \textbf{Representative Test IDs} \\
		\midrule
		1. Authentication \& Account Management & FR 1--4; SR AR 1--4; SR3 & AUTH-UT-001--003 \\
		2. Organization Linking \& Access Codes & FR 1.2--1.3; FR 2--3 & ORG-UT-001--002 \\
		3. Exercise Program Management & FR 5--6; FR 10--11; FR 22; FR 27 & PROG-UT-001--003 \\
		4. Video Capture \& Review & FR 7; FR 15; FR 26 & CAP-UT-001--003 \\
		5. Computer Vision \& Form Analysis & FR 8; FR 12--14; FR 16; PR-PAR 1--2; PR-SCR 1 & CV-UT-001--003 \\
		6. Reporting \& Dashboards & FR 17--20; FR 25 & RPT-UT-001--002 \\
		7. Data Security \& Privacy & FR 21--23; SR-IR 1; SR-PR 1--2 & DATA-UT-001--002 \\
		8. Performance \& Latency & PR-SL 1--2 & PERF-UT-001--002 \\
		9. Reliability \& Offline Operation & PR-RFT 2; MS-AD 1 & REL-UT-001--002 \\
		10. Security \& Backup & SR AR 3; MS-MNT 1 & SEC-UT-001--002 \\
		11. Scalability \& Extensibility & PR-CR 1; PR-SER 1 & EXT-UT-001 \\
		\bottomrule
	\end{tabular}
	\caption{Traceability Between Test Cases and Modules}
\end{table}

This matrix demonstrates coverage of all functional and nonfunctional requirements through at least one representative unit test per module.


\section{Appendix}

Additional project information to expand on the previous sections. Mostly example or sample information at this time.

\subsection{Example script}

\noindent Example script layout for system testing: \\
/* \\
\# Create user account testing \\
New account creation (user\_acc\_create) \\
Invalid format of field inputs (invalid\_formatCreate) \\
Invalid sign up credentials (cannot\_verify\_creds) \\
Presence of an existing account (acc\_exists) \\ \\
\# PT focused testing \\
PT successfully sends invite to patient (PT\_sends\_invite) \\
PT adds/modified patient exercise (PT\_adjusts\_exercise) \\
PT deletes patient account (PT\_deletes\_patient\_acc) \\ \\
\# Patient focused testing \\
System\_exercise\_overview \\
Video\_store \\
Video\_retrieval \\
System\_analysis\_feedback \\
Post\_ex\_patient\_feedback \\
Incorrect\_form\_adjustment \\ \\
\# NFR testing \\
Ex\_int\_speed \\
Fail\_recovery \\
Interrupted\_vid \\
\*/


\subsection{Usability Survey Questions}

\noindent \textbf{Advisor SRS Questions}:
\begin{enumerate}
  \item Are there any sections of the requirements which are found to be confusing or hard to measure?
  \item Are there any sections of the system you feel the requirements do not expand on?
  \item Are there any sections of the system you feel the requirements expand on that is out of scope?
  \item Do you feel the requirements are extensive enough?
  \item Do you feel the requirements are tracked well?
\end{enumerate}

\noindent \textbf{Design Questionnaire Questions}:

“Strongly Disagree” to “Strongly Agree” style questioning.
\begin{enumerate}
  \item The system navigation is intuitive.
  \item The system’s functions are all easily accessible.
  \item The system helped set up your exercises for success.
  \item The system enables you to complete your exercises with better guidance.
  \item The system enables you to complete your exercises with better confidence.
  \item The feedback returned during your exercise was helpful.
  \item The overview data from your exercise was helpful.
  \item You would you use this software for your exercises again.
\end{enumerate}
Could have a final open-ended question based on user pain points with the system I.e. what did not work well for individuals. \\ \\
\noindent \textbf{Software Validation Questionnaire Questions}:

“Strongly Disagree” to “Strongly Agree” style questioning.
\begin{enumerate}
  \item This software assisted you in completing your exercises.
  \item This software was easy to use.
  \item This software was easy to navigate.
  \item The feedback from the system assisted you in fixing your form.
  \item The feedback from the system generated greater confidence in your form.
  \item The feedback from the system was meaningful and helpful.
  \item The data generated from your exercises was helpful.
  \item The system responded how you expected it to during usage.
  \item The system clearly described what was going on.
  \item You felt adequately trained to use the system from the system dialogues.
  \item You felt safe while using the system.
\end{enumerate}

\noindent \textbf{Ui\_testing Survey}:
\begin{enumerate}
  \item How would you rate your overall experience with the application (1 - poor, 5 - excellent)?
  \item How would you rate the visual appeal of the application (1 - poor, 5 - excellent)? 
  \item How easy was it to navigate through the application? (1 - difficult, 5 - easy)
  \begin{enumerate}
    \item Were the tutorials helpful? (1 - not helpful, 5 - very helpful)
    \item How logical was the arrangement of the different components within the system (1 - very illogical, 5 - very logical)?
  \end{enumerate}
  \item Which areas could use some improvement?
  \begin{enumerate}
    \item Navigation 
    \item Menu structure
    \item Colour/contrast scheme
    \item Error messages 
    \item Other (With open text entry)
  \end{enumerate}
\end{enumerate}

\noindent \textbf{Patient Post-Exercise Feedback Form}:
\begin{enumerate}
  \item How did you feel before the exercise? (1 - good, 5-bad)
  \begin{enumerate}
    \item Did you have tightness? Yes/no
    \item Did you have any pain? Yes/no.
    \begin{enumerate}
      \item How bad was your pain? (1 - no pain, 10 - unbearable pain)
      \item What was the perceived level of pain/difficulty? (1 - very easy, 10 - very difficult)
      \item Where was your pain located? (Open text entry)
    \end{enumerate}
    \item Have you noticed any recent changes? (Open text entry)
  \end{enumerate}
  \item If you had any pain DURING the exercise, how bad was it? (1 - no pain, 10 - unbearable pain)
  \item How did you feel after the exercise? (1 - good, 5-bad)
  \begin{enumerate}
    \item Did you have tightness? Yes/no
    \item Did you have any pain? Yes/no.
    \begin{enumerate}
      \item How bad was your pain? (1 - no pain, 10 - unbearable pain)
      \item What was the perceived level of pain/difficulty? (1 - very easy, 10 - very difficult)
      \item Where was your pain located? (Open text entry)
    \end{enumerate}
    \item Have you noticed any recent changes? (Open text entry)
  \end{enumerate}
  \item How challenging was the exercise today? (1 - easy, 5 - very hard)
  \begin{enumerate}
    \item Do you have any feedback for your physiotherapist? (Open text entry)
  \end{enumerate}
  \item In your opinion, is this treatment plan helping your health? Yes/no. Choose the option that describes how the plan is helping (Can choose multiple).
  \begin{enumerate}
    \item My movement
    \item My flexibility
    \item My pain
    \item My mood
  \end{enumerate}
  \item In your opinion, how difficult is the current prescribed exercise plan? (1 - easy, 5 - very difficult)
  \begin{enumerate}
    \item Should your physiotherapist change the difficulty of your plan? (1 - no change, 5 - Large increase in difficulty)
  \end{enumerate}
\end{enumerate}
 
\newpage{}
\bibliographystyle {plainnat}
\bibliography{../../refs/References}   
\newpage{}
\section*{Appendix --- Reflection}

The information in this section will be used to evaluate the team members on the
graduate attribute of Lifelong Learning.

\input{../Reflection.tex}

\begin{enumerate}
  \item What went well while writing this deliverable? 
  \item What pain points did you experience during this deliverable, and how
    did you resolve them?
  \item What knowledge and skills will the team collectively need to acquire to
  successfully complete the verification and validation of your project?
  Examples of possible knowledge and skills include dynamic testing knowledge,
  static testing knowledge, specific tool usage, Valgrind etc.  You should look to
  identify at least one item for each team member.
  \item For each of the knowledge areas and skills identified in the previous
  question, what are at least two approaches to acquiring the knowledge or
  mastering the skill?  Of the identified approaches, which will each team
  member pursue, and why did they make this choice?
\end{enumerate}

\subsection*{Vaisnavi - Reflection}
\begin{enumerate}
  \item I think one main thing that went well while writing this deliverable was our communication as a group and how that played into the division of work. We had two sub-teams where Matthew and I were focusing on the POC, while Cieran, Maham and Eman were focused on this VnV deliverable. Our communication throughout aided in making this process of working on the deliverable seamless as everyone seemed to be very looped in on what each sub-team was working on. We had our weekly meetings to touch base and ask for any help if it was needed, and it worked very well.
  \item I think one of the main pain points we came across was as a result of us working on the POC, it was a little bit hard at the start to fully keep up with updates for the VnV Plan. However, as mentioned above, with sharing our progress and updates during our team meetings, this helped in providing more clarity and consistency between all group members, and it allowed everyone to be on the same page.
  \item \textbf{GROUP Q}
  \item To improve the knowledge of OpenCV and MediaPipe pose detection for accurate testing of body landmark tracing, we could do the two approaches:
  \begin{enumerate}
    \item Follow documentation and existing examples to understand how to properly utilize OpenCV and Mediapipe to accurately detect the varying body landmarks through pose detection.
    \item Try to work with small test cases and work on small files to test it out on different body landmarks and see how it works.
  \end{enumerate}
\end{enumerate}

\subsection*{Maham - Reflection}
\begin{enumerate}
  \item The team was split into two groups: one was responsible for the POC demonstration, and the other was completing the VnV report. All members were in constant communication even though they had other (very) important deadlines to meet. The sub-teams functioned very well, because each team had members who were (somewhat) more comfortable in the area they were responsible for, and others that wanted to expand their skillset. There was a good balance present. In the end, all members delivered work of the highest calibre and I am happy to be a part of such a great team.
  \item I do not think it was a pain point but I had underestimated the amount of thought and importance that my portion would entail. I think it was challenging balancing the high priority of this deliverable (my part) and my other important (health) engagements as well. However, I used all the time management skills that I have learned and effectively managed to chunk my work into more realizable chunks while ensuring I am taking care of my well-being and delivered some good quality work.
  \item \textbf{GROUP Q}
  \item I had previously acquired some knowledge on software testing and have always been fascinated by it. But to really master the application, I had to be very detail oriented. I did a deep dive into the various components of our system. However, I realized that I had to be careful and not make my focus too narrow - because I was developing system tests. I had to look at the overall picture. I did this in 2 ways. Firstly, I researched ‘system testing’ to further increase my knowledge and get a more concrete idea. However, since I had prior knowledge, I found most of the resources online just broke down system testing into different unit tests. I did not want to pursue that because Eman was working on the unit testing component. I spoke to my teammates on the possible overlap between unit and system testing, and was more confident in my (potential) analysis and approach. In the end, I refined my testing knowledge and by designing different test cases for any possible scenario, and making sure the system was approached from various aspects without being too focused on individual modules. 
\end{enumerate}

\subsection*{Matthew - Reflection}
\begin{enumerate}
  \item I think one thing that went well when writing this deliverable was how we divided these two weeks up. We had taken the advice from the professor to begin development on the POC while working on the VnV plan. The team was divided into myself and Vais on the POC while Cieran, Maham and Eman were working on the VnV plan. This allowed us to do brief check-ins with each other while beginning development and thus more time to develop before the POC demo.
  \item I think some pain points during this deliverable was tied to external workloads. Since half of the deliverable encompassed reading week, this meant that the team had midterms and assignments for other classes due before this deliverable and managing that was more difficult. We resolved this by being transparent with each of our workloads, by letting each other know in our Teams chat if we were going to be busy and when the next time we would be available is. This allowed us to plan the deliverable better and achieve internal deadlines.
  \item \textbf{GROUP Q}
  \item I would like to improve my knowledge of stakeholder understanding. I think that with stakeholder understanding it is important as they are the target audience for the application and their needs are to be met even if it differs from our own since they are the target. My approach for that will consist of understanding interviews as well as conversations with the stakeholder during testing. Another approach will be different user groups to see if our initial persona and ideas for the stakeholder would make sense. \\
  I think the is important since it is easy to develop requirements by making assumptions but to bridge the gap to the target audience is crucial to make a product or service that they would actually use.

\end{enumerate}

\subsection*{Cieran - Reflection}
\begin{enumerate}
  \item Team communication as always but I won’t expand further on that. There was a decision to split the team into two subteams; a proof of concept team to begin the implementation and a verification and validation team to tackle this document. I feel the subteams have both been incredibly efficient in their work and the communication between the two has not faltered. The proof of concept team worked well together in overall and in subteam meetings to grasp an understanding of the software packages we will be utilizing while the verification and validation team ensured the entire system was covered within the tests outlined in this document.
  \item I expanded on it during the last deliverable, that each person had this ‘own version’ of the system in their mind and we were all building towards a different thing when the assignment started. That wasn’t entirely gone here and we hit a few bumps worrying about what each of the others was including within the system while thinking we don’t need that to be included. Through team meetings and a general walkthrough of the flow of the program, we all got on the same page and managed to generate tests and feel as if we’ve properly outlined a system to ensure the system is designed properly.
  \item \textbf{GROUP Q}
  \item I feel my personal interaction with test suite development lacks for the requirements of this project. To advance this I will attempt to do the following:
  \begin{enumerate}
    \item Follow online guides, walkthroughs, and documentation on the software packages (Python Extension, pytest, pylint, etc.) to ensure the software testing suite developed contains both meaningful tests derived from this document while also covering all the software.
    \item Generate and execute tests in advancing complexity to ensure an understanding of the routines gone through with a testing suite. ‘Getting my hands dirty’.
  \end{enumerate}
\end{enumerate}

\subsection*{Eman - Reflection}
    \begin{enumerate}
      \item The team worked together effectively by dividing into two focused groups, one for the POC demonstration and another for the VnV plan. Communication was strong throughout, and everyone contributed their strengths while helping others learn. The collaboration created a good balance where team members could build on their existing skills while developing new ones.
      \item My main challenge was ensuring complete test coverage across all system modules. I initially underestimated how much detail was needed to properly test each component in isolation. Identifying all the edge cases and understanding how to structure effective unit tests took more effort than expected. I resolved this by creating a systematic checklist that mapped each module to its requirements, then worked through them methodically. Collaborating with teammates to review my test cases also helped catch gaps I had missed.
      \item \textbf{GROUP Q}
      \item I focused on mastering unit testing and test coverage. My first approach was working through practical examples to understand fixtures, mocking, and test isolation. My second approach was carefully reviewing our system requirements and architecture to understand what each module should do. I chose this combination because I needed both the technical testing skills and domain knowledge about our system. Understanding the requirements helped me design tests that actually validated meaningful functionality rather than just checking basic operations. This dual approach allowed me to create comprehensive unit tests covering all eleven modules in our system.
    \end{enumerate}


\end{document}